\subsection{Module Giỏ hàng và Theo dõi đơn hàng}
\label{sec:cart-tracking}

\paragraph{Giỏ hàng (\texttt{useCartStore}).}
Giỏ hàng được quản lý bằng Zustand và đồng bộ với \texttt{localStorage} (khóa \texttt{brewlite\_cart}) nên không mất dữ liệu khi tải lại trang, dùng chung với giỏ hàng ở header và trang thanh toán
Mỗi dòng được định danh bằng khoá \texttt{productId|size|toppings}, vì vậy cùng một đồ uống với cùng
size và topping sẽ được gộp số lượng. Giá trong giỏ chỉ dùng để hiển thị; tổng tiền thực tế
luôn do Backend tính lại.

\paragraph{API theo dõi đơn.}
\begin{itemize}
  \item \texttt{GET /orders/me}: trả về các đơn của người dùng đăng nhập, mới nhất trước, kèm sản phẩm và thanh toán.
  \item \texttt{PATCH /orders/:id/status}: đổi trạng thái theo máy trạng thái
        \texttt{PENDING $\rightarrow$ PAID $\rightarrow$ COMPLETED}, có thể huỷ (\texttt{CANCELLED}) khi đơn đang
        \texttt{PENDING} hoặc \texttt{PAID}. Chuyển sai thứ tự trả về \texttt{400}; khách hàng chỉ được huỷ đơn của chính mình,
        các thay đổi khác thuộc về nhân viên (role \texttt{ADMIN}/\texttt{STAFF}).
\end{itemize}

\paragraph{Order Tracking Stepper.}
Trang \texttt{/orders} hiển thị tiến trình theo bảng ánh xạ sau, tự làm mới mỗi 5 giây và có đủ các trạng thái Loading, Empty, Error.

\begin{table}[H]
  \centering
  \begin{tabular}{|l|l|}
    \hline
    \textbf{Trạng thái trong CSDL} & \textbf{Hiển thị trên Stepper} \\ \hline
    \texttt{PENDING}   & Đã đặt \\ \hline
    \texttt{PAID}      & Đang pha chế \\ \hline
    \texttt{COMPLETED} & Sẵn sàng nhận \\ \hline
    \texttt{CANCELLED} & Đã huỷ \\ \hline
  \end{tabular}
  \caption{Ánh xạ \texttt{OrderStatus} sang các bước của Stepper}
\end{table}
